\section{Discussion}\label{sec:discussion}

In this paper we have introduced a novel statistical test to assess treatment effect on continuous outcomes where one value has special meaning (e.g., all deceased are assigned lowest possible value). The procedure is potentially much more powerful than the current best practice which is to use Wilcoxon-type tests. The proposed method includes both a fully parametric approach and a semi-parametric approach where the latter makes no assumptions on the distribution of the continuous part of the outcome. In all settings this new method not only provides an effect measure but also effect parameters with associated confidence intervals. The test is implemented in an R package available on GitHub.

The simulation study and in particular the real data example demonstrate the efficiency and power gain. It is noted that unlike the Wilcoxon test, our proposed method can easily be extended to include covariates \citep{owen1991empirical}. It is therefore not only useful in RCT settings but also in epidemiological studies.

\textbf{[TODO: Do we need something about that this is very useful when needing to prespecify a test in a RCT?]}

\textbf{[TODO: check verb tense throughout]}
